% lab/latex/preamble-ja.tex
% 日本語主体・欧文混在の学術論文用 共通プリアンブル（LuaLaTeX + jlreq）
% 使い方: % lab/latex/preamble-ja.tex
% 日本語主体・欧文混在の学術論文用 共通プリアンブル（LuaLaTeX + jlreq）
% 使い方: % lab/latex/preamble-ja.tex
% 日本語主体・欧文混在の学術論文用 共通プリアンブル（LuaLaTeX + jlreq）
% 使い方: % lab/latex/preamble-ja.tex
% 日本語主体・欧文混在の学術論文用 共通プリアンブル（LuaLaTeX + jlreq）
% 使い方: \input{../../lab/latex/preamble-ja}  ← 各プロジェクトの main.tex から

%%% フォント %%%
\usepackage{luatexja-fontspec}
\usepackage{iftex}
% 和文明朝: 原ノ味明朝（TeX Live同梱、全OS対応）
\setmainjfont{HaranoAjiMincho}
% 和文ゴシック: OS別にフォールバック
%   macOS: ヒラギノ角ゴシック ProN
%   Linux/Windows: Noto Sans CJK JP → 原ノ味ゴシック（TeX Live同梱）
\IfFontExistsTF{Hiragino Kaku Gothic ProN}{%
  \setsansjfont{Hiragino Kaku Gothic ProN}[
    UprightFont = {Hiragino Kaku Gothic ProN W3},
    BoldFont    = {Hiragino Kaku Gothic ProN W6},
  ]%
}{%
  \IfFontExistsTF{Noto Sans CJK JP}{%
    \setsansjfont{Noto Sans CJK JP}%
  }{%
    \setsansjfont{HaranoAjiGothic}%
  }%
}
% 見出しをゴシック体にする
\AddToHook{begindocument/before}{%
  \ModifyHeading{chapter}{font={\sffamily\bfseries}}%
  \ModifyHeading{section}{font={\sffamily\bfseries}}%
  \ModifyHeading{subsection}{font={\sffamily\bfseries}}%
}

% 部の見出し: 番号と題名を別行に、行間を広げる
\makeatletter
\renewcommand{\part}{%
  \clearpage
  \thispagestyle{empty}%
  \null\vfil
  \secdef\@part\@spart}
\def\@part[#1]#2{%
  \refstepcounter{part}%
  \addcontentsline{toc}{part}{\thepart\hspace{1em}#1}%
  \begin{center}
    {\huge\sffamily\bfseries 第\thepart 部}\\[1.5em]
    {\LARGE\sffamily\bfseries #2}
  \end{center}
  \vfil}
\def\@spart#1{%
  \begin{center}
    {\LARGE\sffamily\bfseries #1}
  \end{center}
  \vfil}
\makeatother
% 欧文: TeX Gyre Termes（Times系、TeX Live同梱）
\setmainfont{TeX Gyre Termes}

%%% 脚注番号の括弧を外す %%%
\renewcommand{\thefootnote}{\arabic{footnote}}

%%% 節見出しの番号を「第n節」にする %%%
\renewcommand{\thesection}{第\arabic{section}節}
\renewcommand{\thesubsection}{\arabic{section}.\arabic{subsection}}

%%% 目次の体裁 %%%
% 部: 章と同じインデント体裁で、太字ゴシック
\makeatletter
\renewcommand*{\l@part}[2]{%
  \ifnum \c@tocdepth >-2\relax
    \addpenalty{-\@highpenalty}%
    \addvspace{1em \@plus\p@}%
    \@dottedtocline{-1}{0em}{5em}{\bfseries\sffamily #1}{\bfseries #2}%
    \nobreak
  \fi}
\makeatother

% oneside: 空白ページの挿入を防止
\let\cleardoublepage\clearpage

%%% 参考文献 %%%
\usepackage[
  backend=biber,
  style=verbose-note,    % 脚注内にフル書誌→短縮形（人文系標準）
  autocite=footnote,
  defernumbers=true,     % 分割出力時の番号管理
  sorting=nyt,
  abbreviate=false,
  language=auto,         % langid による言語別自動切替
]{biblatex}

% 参考文献一覧の見出し階層定義
\defbibheading{subbibliography}[\bibname]{%
  \section*{#1}}
\defbibheading{subsubbibliography}[\bibname]{%
  \subsection*{#1}}

%%% 書誌フォーマット（修論準拠）%%%
%
% langid による言語別自動切替:
%   french  → « guillemets », "dans", "et", "éd."
%   english → "quotation marks", "in", "and", "ed."
%   japanese → キーワードで別処理
% .bibエントリに langid = {french/english/japanese} を指定すること。

% 脚注・参考文献一覧ともに姓-名の順に統一
\DeclareNameAlias{sortname}{family-given}
\DeclareNameAlias{default}{family-given}

% 欧語: 姓をスモールキャピタルにする
% \ifkeyword は短縮引用で効かないため、LuaLaTeX で文字コード判定
\directlua{
  function is_cjk(s)
    local first = utf8.codepoint(s, 1)
    if first >= 0x3000 and first <= 0x9FFF then return true end
    if first >= 0xF900 and first <= 0xFAFF then return true end
    return false
  end
}
\renewcommand*{\mkbibnamefamily}[1]{%
  \directlua{
    if is_cjk("\luaescapestring{#1}") then
      tex.sprint("\luaescapestring{#1}")
    else
      tex.sprint("\\textsc{\luaescapestring{#1}}")
    end
  }%
}

% 日本語文献: 姓名スペース除去 + 句読点を日本語に
% 日本語用の句読点設定マクロ
\newcommand*{\setjapanesepunct}{%
  \renewcommand*{\revsdnamepunct}{}%
  \renewcommand*{\bibnamedelimd}{}%
  \renewcommand*{\newunitpunct}{\addspace}%
  \renewcommand*{\labelnamepunct}{\addspace}%
  \renewcommand*{\subtitlepunct}{：}%
  \renewcommand*{\bibpagespunct}{、}%
  \renewcommand*{\multicitedelim}{、}%
  \renewcommand*{\finentrypunct}{。}%
  \DeclareDelimFormat{nameyeardelim}{、}%
  \DeclareDelimFormat{nonameyeardelim}{、}%
  \DeclareDelimFormat[bib,biblist]{multinamedelim}{・}%
  \DeclareDelimFormat[bib,biblist]{finalnamedelim}{・}%
  \DeclareDelimFormat{titleyeardelim}{、}%
  \DeclareDelimFormat{publocpunct}{、}%
  \DeclareDelimFormat{publocdelim}{、}%
  \DeclareDelimFormat{publisherdelim}{、}%
  \renewcommand*{\newblockpunct}{\addspace}%
  \DeclareDelimFormat{postnotedelim}{、}%
}
% 欧語用の句読点設定マクロ
\newcommand*{\setwesternpunct}{%
  \renewcommand*{\revsdnamepunct}{\addcomma\space}%
  \renewcommand*{\bibnamedelimd}{\addspace}%
  \renewcommand*{\newunitpunct}{\adddot\space}%
  \renewcommand*{\subtitlepunct}{\addcolon\space}%
  \renewcommand*{\finentrypunct}{\adddot}%
}
% 日本語: 出版年の後に「年」を付ける
\DeclareFieldFormat{datetext}{%
  \ifkeyword{japanese}{#1年}{#1}}
\renewbibmacro*{date}{%
  \printtext[datetext]{\printdate}%
}

\AtEveryBibitem{\ifkeyword{japanese}{\setjapanesepunct}{\setwesternpunct}}
\AtEveryCitekey{\ifkeyword{japanese}{\setjapanesepunct}{\setwesternpunct}}

% 脚注末尾: 日本語「。」/ 欧語「.」
\renewcommand{\bibfootnotewrapper}[1]{\bibsentence #1}

% 短縮引用: 日本語では「著者、前掲書」形式にする
\renewbibmacro*{cite:short}{%
  \ifkeyword{japanese}%
    {\printnames{labelname}、前掲書}%
    {\printnames{labelname}%
     \setunit*{\printdelim{nametitledelim}}%
     \printtext[bibhyperref]{%
       \printfield[citetitle]{labeltitle}}}%
}

% autocite: 末尾句点を postnote の後に出す
\DeclareCiteCommand{\autocite}[\mkbibfootnote]
  {\usebibmacro{prenote}}
  {\usebibmacro{citeindex}%
   \usebibmacro{cite}}
  {\multicitedelim}
  {\usebibmacro{cite:postnote}\finentrypunct}

% 出版社と年の間の区切りを日本語では「、」に
\renewbibmacro*{publisher+location+date}{%
  \printlist{location}%
  \iflistundef{publisher}%
    {\setunit*{\addcomma\space}}%
    {\setunit*{\ifkeyword{japanese}{、}{\addcomma\space}}}%
  \printlist{publisher}%
  \setunit*{\ifkeyword{japanese}{、}{\addcomma\space}}%
  \usebibmacro{date}%
  \newunit%
}

% 論文・章題の引用符
% 日本語「」、フランス語 « »、英語 " "
\newcommand*{\langquote}[1]{%
  \iffieldequalstr{langid}{french}%
    {\guillemotleft\addnbspace #1\addnbspace\guillemotright}%
    {\mkbibquote{#1}}%
}
\DeclareFieldFormat[article]{title}{%
  \ifkeyword{japanese}{「#1」}{\langquote{#1}}}
\DeclareFieldFormat[incollection]{title}{%
  \ifkeyword{japanese}{「#1」}{\langquote{#1}}}
\DeclareFieldFormat[inproceedings]{title}{%
  \ifkeyword{japanese}{「#1」}{\langquote{#1}}}
\DeclareFieldFormat[thesis]{title}{%
  \ifkeyword{japanese}{「#1」}{\langquote{#1}}}

% 日本語の書名を『』で囲む（イタリック→ゴシック化を防止）
\DeclareFieldFormat[book]{title}{%
  \ifkeyword{japanese}{『#1』}{\mkbibemph{#1}}}
\DeclareFieldFormat[book]{subtitle}{%
  \ifkeyword{japanese}{#1}{\mkbibemph{#1}}}
\DeclareFieldFormat[incollection]{booktitle}{%
  \ifkeyword{japanese}{『#1』}{\mkbibemph{#1}}}
\DeclareFieldFormat[book]{citetitle}{%
  \ifkeyword{japanese}{『#1』}{\mkbibemph{#1}}}
\DeclareFieldFormat[incollection]{citetitle}{%
  \ifkeyword{japanese}{「#1」}{\langquote{#1}}}
\DeclareFieldFormat[article]{citetitle}{%
  \ifkeyword{japanese}{「#1」}{\langquote{#1}}}

% 日本語: "In:" を抑制、欧語: 大文字 "In:" を維持
% 日本語: 抑制、フランス語: dans、英語: in
\renewbibmacro*{in:}{%
  \ifkeyword{japanese}%
    {}%
    {\iffieldequalstr{langid}{french}%
      {\printtext{dans\intitlepunct}}%
      {\printtext{in\intitlepunct}}}%
}

% 日本語: 短縮引用「前掲注X」+ pages 表記
\DefineBibliographyStrings{english}{%
  seenote = {前掲注},
}
% 日本語文献の pages を日本語化
\DeclareFieldFormat[incollection]{pages}{%
  \ifkeyword{japanese}{#1頁}{pp.\addnbspace #1}}
\DeclareFieldFormat[article]{pages}{%
  \ifkeyword{japanese}{#1頁}{pp.\addnbspace #1}}

% 日本語: "Edited by ..." を抑制、欧語: 名前 (eds.) 形式
\renewbibmacro*{byeditor+others}{%
  \ifkeyword{japanese}%
    {}%
    {\ifnameundef{editor}%
      {}%
      {\printnames[byeditor]{editor}%
       \addspace\printeditorlabel
       \clearname{editor}%
       \newunit}%
    }%
}

% 編者の括弧付き表記ヘルパー
\newcommand*{\printeditorlabel}{%
  \iffieldequalstr{langid}{french}%
    {\printtext[parens]{\ifnumgreater{\value{editor}}{1}{éds\adddot}{éd\adddot}}}%
    {\printtext[parens]{\ifnumgreater{\value{editor}}{1}{eds\adddot}{ed\adddot}}}%
}

% editor+othersstrg / editor+others: 括弧付き (eds.)/(éds.) を出力
\renewbibmacro*{editor+othersstrg}{\printeditorlabel}
\renewbibmacro*{editor+others}{%
  \ifboolexpr{
    test \ifuseeditor
    and
    not test {\ifnameundef{editor}}
  }
    {\printnames{editor}%
     \addspace\printeditorlabel
     \clearname{editor}}
    {}%
}

% book型: 編者名 (éds.)/(eds.) の形式
\renewbibmacro*{bbx:editor}{%
  \ifnameundef{editor}{}%
    {\printnames{editor}%
     \usebibmacro{editor+othersstrg}%
     \clearname{editor}%
     \setunit{\labelnamepunct}}%
}

% フランス語文献: "and" → "et"
\AtEveryBibitem{%
  \iffieldequalstr{langid}{french}%
    {\DeclareDelimFormat{finalnamedelim}{\addspace et\addspace}}%
    {}%
}
\AtEveryCitekey{%
  \iffieldequalstr{langid}{french}%
    {\DeclareDelimFormat{finalnamedelim}{\addspace et\addspace}}%
    {}%
}

%%% 組版 %%%
\usepackage{setspace}
\onehalfspacing
\usepackage[
  top=30mm,bottom=30mm,
  left=30mm,right=25mm,
]{geometry}


%%% 引用符 %%%
\usepackage[autostyle=true, french=guillemets, english=american]{csquotes}

%%% ハイパーリンク %%%
\usepackage[dvipsnames,svgnames,x11names]{xcolor}
\usepackage{hyperref}
\hypersetup{
  colorlinks=true,
  linkcolor=black,
  citecolor=black,
  urlcolor=blue!60!black,
  bookmarksnumbered=true,
}

%%% その他 %%%
\usepackage{graphicx}
\usepackage{booktabs}
\usepackage{wrapfig}

%%% 見出し %%%
\setcounter{secnumdepth}{1}
\setcounter{tocdepth}{2}
  ← 各プロジェクトの main.tex から

%%% フォント %%%
\usepackage{luatexja-fontspec}
\usepackage{iftex}
% 和文明朝: 原ノ味明朝（TeX Live同梱、全OS対応）
\setmainjfont{HaranoAjiMincho}
% 和文ゴシック: OS別にフォールバック
%   macOS: ヒラギノ角ゴシック ProN
%   Linux/Windows: Noto Sans CJK JP → 原ノ味ゴシック（TeX Live同梱）
\IfFontExistsTF{Hiragino Kaku Gothic ProN}{%
  \setsansjfont{Hiragino Kaku Gothic ProN}[
    UprightFont = {Hiragino Kaku Gothic ProN W3},
    BoldFont    = {Hiragino Kaku Gothic ProN W6},
  ]%
}{%
  \IfFontExistsTF{Noto Sans CJK JP}{%
    \setsansjfont{Noto Sans CJK JP}%
  }{%
    \setsansjfont{HaranoAjiGothic}%
  }%
}
% 見出しをゴシック体にする
\AddToHook{begindocument/before}{%
  \ModifyHeading{chapter}{font={\sffamily\bfseries}}%
  \ModifyHeading{section}{font={\sffamily\bfseries}}%
  \ModifyHeading{subsection}{font={\sffamily\bfseries}}%
}

% 部の見出し: 番号と題名を別行に、行間を広げる
\makeatletter
\renewcommand{\part}{%
  \clearpage
  \thispagestyle{empty}%
  \null\vfil
  \secdef\@part\@spart}
\def\@part[#1]#2{%
  \refstepcounter{part}%
  \addcontentsline{toc}{part}{\thepart\hspace{1em}#1}%
  \begin{center}
    {\huge\sffamily\bfseries 第\thepart 部}\\[1.5em]
    {\LARGE\sffamily\bfseries #2}
  \end{center}
  \vfil}
\def\@spart#1{%
  \begin{center}
    {\LARGE\sffamily\bfseries #1}
  \end{center}
  \vfil}
\makeatother
% 欧文: TeX Gyre Termes（Times系、TeX Live同梱）
\setmainfont{TeX Gyre Termes}

%%% 脚注番号の括弧を外す %%%
\renewcommand{\thefootnote}{\arabic{footnote}}

%%% 節見出しの番号を「第n節」にする %%%
\renewcommand{\thesection}{第\arabic{section}節}
\renewcommand{\thesubsection}{\arabic{section}.\arabic{subsection}}

%%% 目次の体裁 %%%
% 部: 章と同じインデント体裁で、太字ゴシック
\makeatletter
\renewcommand*{\l@part}[2]{%
  \ifnum \c@tocdepth >-2\relax
    \addpenalty{-\@highpenalty}%
    \addvspace{1em \@plus\p@}%
    \@dottedtocline{-1}{0em}{5em}{\bfseries\sffamily #1}{\bfseries #2}%
    \nobreak
  \fi}
\makeatother

% oneside: 空白ページの挿入を防止
\let\cleardoublepage\clearpage

%%% 参考文献 %%%
\usepackage[
  backend=biber,
  style=verbose-note,    % 脚注内にフル書誌→短縮形（人文系標準）
  autocite=footnote,
  defernumbers=true,     % 分割出力時の番号管理
  sorting=nyt,
  abbreviate=false,
  language=auto,         % langid による言語別自動切替
]{biblatex}

% 参考文献一覧の見出し階層定義
\defbibheading{subbibliography}[\bibname]{%
  \section*{#1}}
\defbibheading{subsubbibliography}[\bibname]{%
  \subsection*{#1}}

%%% 書誌フォーマット（修論準拠）%%%
%
% langid による言語別自動切替:
%   french  → « guillemets », "dans", "et", "éd."
%   english → "quotation marks", "in", "and", "ed."
%   japanese → キーワードで別処理
% .bibエントリに langid = {french/english/japanese} を指定すること。

% 脚注・参考文献一覧ともに姓-名の順に統一
\DeclareNameAlias{sortname}{family-given}
\DeclareNameAlias{default}{family-given}

% 欧語: 姓をスモールキャピタルにする
% \ifkeyword は短縮引用で効かないため、LuaLaTeX で文字コード判定
\directlua{
  function is_cjk(s)
    local first = utf8.codepoint(s, 1)
    if first >= 0x3000 and first <= 0x9FFF then return true end
    if first >= 0xF900 and first <= 0xFAFF then return true end
    return false
  end
}
\renewcommand*{\mkbibnamefamily}[1]{%
  \directlua{
    if is_cjk("\luaescapestring{#1}") then
      tex.sprint("\luaescapestring{#1}")
    else
      tex.sprint("\\textsc{\luaescapestring{#1}}")
    end
  }%
}

% 日本語文献: 姓名スペース除去 + 句読点を日本語に
% 日本語用の句読点設定マクロ
\newcommand*{\setjapanesepunct}{%
  \renewcommand*{\revsdnamepunct}{}%
  \renewcommand*{\bibnamedelimd}{}%
  \renewcommand*{\newunitpunct}{\addspace}%
  \renewcommand*{\labelnamepunct}{\addspace}%
  \renewcommand*{\subtitlepunct}{：}%
  \renewcommand*{\bibpagespunct}{、}%
  \renewcommand*{\multicitedelim}{、}%
  \renewcommand*{\finentrypunct}{。}%
  \DeclareDelimFormat{nameyeardelim}{、}%
  \DeclareDelimFormat{nonameyeardelim}{、}%
  \DeclareDelimFormat[bib,biblist]{multinamedelim}{・}%
  \DeclareDelimFormat[bib,biblist]{finalnamedelim}{・}%
  \DeclareDelimFormat{titleyeardelim}{、}%
  \DeclareDelimFormat{publocpunct}{、}%
  \DeclareDelimFormat{publocdelim}{、}%
  \DeclareDelimFormat{publisherdelim}{、}%
  \renewcommand*{\newblockpunct}{\addspace}%
  \DeclareDelimFormat{postnotedelim}{、}%
}
% 欧語用の句読点設定マクロ
\newcommand*{\setwesternpunct}{%
  \renewcommand*{\revsdnamepunct}{\addcomma\space}%
  \renewcommand*{\bibnamedelimd}{\addspace}%
  \renewcommand*{\newunitpunct}{\adddot\space}%
  \renewcommand*{\subtitlepunct}{\addcolon\space}%
  \renewcommand*{\finentrypunct}{\adddot}%
}
% 日本語: 出版年の後に「年」を付ける
\DeclareFieldFormat{datetext}{%
  \ifkeyword{japanese}{#1年}{#1}}
\renewbibmacro*{date}{%
  \printtext[datetext]{\printdate}%
}

\AtEveryBibitem{\ifkeyword{japanese}{\setjapanesepunct}{\setwesternpunct}}
\AtEveryCitekey{\ifkeyword{japanese}{\setjapanesepunct}{\setwesternpunct}}

% 脚注末尾: 日本語「。」/ 欧語「.」
\renewcommand{\bibfootnotewrapper}[1]{\bibsentence #1}

% 短縮引用: 日本語では「著者、前掲書」形式にする
\renewbibmacro*{cite:short}{%
  \ifkeyword{japanese}%
    {\printnames{labelname}、前掲書}%
    {\printnames{labelname}%
     \setunit*{\printdelim{nametitledelim}}%
     \printtext[bibhyperref]{%
       \printfield[citetitle]{labeltitle}}}%
}

% autocite: 末尾句点を postnote の後に出す
\DeclareCiteCommand{\autocite}[\mkbibfootnote]
  {\usebibmacro{prenote}}
  {\usebibmacro{citeindex}%
   \usebibmacro{cite}}
  {\multicitedelim}
  {\usebibmacro{cite:postnote}\finentrypunct}

% 出版社と年の間の区切りを日本語では「、」に
\renewbibmacro*{publisher+location+date}{%
  \printlist{location}%
  \iflistundef{publisher}%
    {\setunit*{\addcomma\space}}%
    {\setunit*{\ifkeyword{japanese}{、}{\addcomma\space}}}%
  \printlist{publisher}%
  \setunit*{\ifkeyword{japanese}{、}{\addcomma\space}}%
  \usebibmacro{date}%
  \newunit%
}

% 論文・章題の引用符
% 日本語「」、フランス語 « »、英語 " "
\newcommand*{\langquote}[1]{%
  \iffieldequalstr{langid}{french}%
    {\guillemotleft\addnbspace #1\addnbspace\guillemotright}%
    {\mkbibquote{#1}}%
}
\DeclareFieldFormat[article]{title}{%
  \ifkeyword{japanese}{「#1」}{\langquote{#1}}}
\DeclareFieldFormat[incollection]{title}{%
  \ifkeyword{japanese}{「#1」}{\langquote{#1}}}
\DeclareFieldFormat[inproceedings]{title}{%
  \ifkeyword{japanese}{「#1」}{\langquote{#1}}}
\DeclareFieldFormat[thesis]{title}{%
  \ifkeyword{japanese}{「#1」}{\langquote{#1}}}

% 日本語の書名を『』で囲む（イタリック→ゴシック化を防止）
\DeclareFieldFormat[book]{title}{%
  \ifkeyword{japanese}{『#1』}{\mkbibemph{#1}}}
\DeclareFieldFormat[book]{subtitle}{%
  \ifkeyword{japanese}{#1}{\mkbibemph{#1}}}
\DeclareFieldFormat[incollection]{booktitle}{%
  \ifkeyword{japanese}{『#1』}{\mkbibemph{#1}}}
\DeclareFieldFormat[book]{citetitle}{%
  \ifkeyword{japanese}{『#1』}{\mkbibemph{#1}}}
\DeclareFieldFormat[incollection]{citetitle}{%
  \ifkeyword{japanese}{「#1」}{\langquote{#1}}}
\DeclareFieldFormat[article]{citetitle}{%
  \ifkeyword{japanese}{「#1」}{\langquote{#1}}}

% 日本語: "In:" を抑制、欧語: 大文字 "In:" を維持
% 日本語: 抑制、フランス語: dans、英語: in
\renewbibmacro*{in:}{%
  \ifkeyword{japanese}%
    {}%
    {\iffieldequalstr{langid}{french}%
      {\printtext{dans\intitlepunct}}%
      {\printtext{in\intitlepunct}}}%
}

% 日本語: 短縮引用「前掲注X」+ pages 表記
\DefineBibliographyStrings{english}{%
  seenote = {前掲注},
}
% 日本語文献の pages を日本語化
\DeclareFieldFormat[incollection]{pages}{%
  \ifkeyword{japanese}{#1頁}{pp.\addnbspace #1}}
\DeclareFieldFormat[article]{pages}{%
  \ifkeyword{japanese}{#1頁}{pp.\addnbspace #1}}

% 日本語: "Edited by ..." を抑制、欧語: 名前 (eds.) 形式
\renewbibmacro*{byeditor+others}{%
  \ifkeyword{japanese}%
    {}%
    {\ifnameundef{editor}%
      {}%
      {\printnames[byeditor]{editor}%
       \addspace\printeditorlabel
       \clearname{editor}%
       \newunit}%
    }%
}

% 編者の括弧付き表記ヘルパー
\newcommand*{\printeditorlabel}{%
  \iffieldequalstr{langid}{french}%
    {\printtext[parens]{\ifnumgreater{\value{editor}}{1}{éds\adddot}{éd\adddot}}}%
    {\printtext[parens]{\ifnumgreater{\value{editor}}{1}{eds\adddot}{ed\adddot}}}%
}

% editor+othersstrg / editor+others: 括弧付き (eds.)/(éds.) を出力
\renewbibmacro*{editor+othersstrg}{\printeditorlabel}
\renewbibmacro*{editor+others}{%
  \ifboolexpr{
    test \ifuseeditor
    and
    not test {\ifnameundef{editor}}
  }
    {\printnames{editor}%
     \addspace\printeditorlabel
     \clearname{editor}}
    {}%
}

% book型: 編者名 (éds.)/(eds.) の形式
\renewbibmacro*{bbx:editor}{%
  \ifnameundef{editor}{}%
    {\printnames{editor}%
     \usebibmacro{editor+othersstrg}%
     \clearname{editor}%
     \setunit{\labelnamepunct}}%
}

% フランス語文献: "and" → "et"
\AtEveryBibitem{%
  \iffieldequalstr{langid}{french}%
    {\DeclareDelimFormat{finalnamedelim}{\addspace et\addspace}}%
    {}%
}
\AtEveryCitekey{%
  \iffieldequalstr{langid}{french}%
    {\DeclareDelimFormat{finalnamedelim}{\addspace et\addspace}}%
    {}%
}

%%% 組版 %%%
\usepackage{setspace}
\onehalfspacing
\usepackage[
  top=30mm,bottom=30mm,
  left=30mm,right=25mm,
]{geometry}


%%% 引用符 %%%
\usepackage[autostyle=true, french=guillemets, english=american]{csquotes}

%%% ハイパーリンク %%%
\usepackage[dvipsnames,svgnames,x11names]{xcolor}
\usepackage{hyperref}
\hypersetup{
  colorlinks=true,
  linkcolor=black,
  citecolor=black,
  urlcolor=blue!60!black,
  bookmarksnumbered=true,
}

%%% その他 %%%
\usepackage{graphicx}
\usepackage{booktabs}
\usepackage{wrapfig}

%%% 見出し %%%
\setcounter{secnumdepth}{1}
\setcounter{tocdepth}{2}
  ← 各プロジェクトの main.tex から

%%% フォント %%%
\usepackage{luatexja-fontspec}
\usepackage{iftex}
% 和文明朝: 原ノ味明朝（TeX Live同梱、全OS対応）
\setmainjfont{HaranoAjiMincho}
% 和文ゴシック: OS別にフォールバック
%   macOS: ヒラギノ角ゴシック ProN
%   Linux/Windows: Noto Sans CJK JP → 原ノ味ゴシック（TeX Live同梱）
\IfFontExistsTF{Hiragino Kaku Gothic ProN}{%
  \setsansjfont{Hiragino Kaku Gothic ProN}[
    UprightFont = {Hiragino Kaku Gothic ProN W3},
    BoldFont    = {Hiragino Kaku Gothic ProN W6},
  ]%
}{%
  \IfFontExistsTF{Noto Sans CJK JP}{%
    \setsansjfont{Noto Sans CJK JP}%
  }{%
    \setsansjfont{HaranoAjiGothic}%
  }%
}
% 見出しをゴシック体にする
\AddToHook{begindocument/before}{%
  \ModifyHeading{chapter}{font={\sffamily\bfseries}}%
  \ModifyHeading{section}{font={\sffamily\bfseries}}%
  \ModifyHeading{subsection}{font={\sffamily\bfseries}}%
}

% 部の見出し: 番号と題名を別行に、行間を広げる
\makeatletter
\renewcommand{\part}{%
  \clearpage
  \thispagestyle{empty}%
  \null\vfil
  \secdef\@part\@spart}
\def\@part[#1]#2{%
  \refstepcounter{part}%
  \addcontentsline{toc}{part}{\thepart\hspace{1em}#1}%
  \begin{center}
    {\huge\sffamily\bfseries 第\thepart 部}\\[1.5em]
    {\LARGE\sffamily\bfseries #2}
  \end{center}
  \vfil}
\def\@spart#1{%
  \begin{center}
    {\LARGE\sffamily\bfseries #1}
  \end{center}
  \vfil}
\makeatother
% 欧文: TeX Gyre Termes（Times系、TeX Live同梱）
\setmainfont{TeX Gyre Termes}

%%% 脚注番号の括弧を外す %%%
\renewcommand{\thefootnote}{\arabic{footnote}}

%%% 節見出しの番号を「第n節」にする %%%
\renewcommand{\thesection}{第\arabic{section}節}
\renewcommand{\thesubsection}{\arabic{section}.\arabic{subsection}}

%%% 目次の体裁 %%%
% 部: 章と同じインデント体裁で、太字ゴシック
\makeatletter
\renewcommand*{\l@part}[2]{%
  \ifnum \c@tocdepth >-2\relax
    \addpenalty{-\@highpenalty}%
    \addvspace{1em \@plus\p@}%
    \@dottedtocline{-1}{0em}{5em}{\bfseries\sffamily #1}{\bfseries #2}%
    \nobreak
  \fi}
\makeatother

% oneside: 空白ページの挿入を防止
\let\cleardoublepage\clearpage

%%% 参考文献 %%%
\usepackage[
  backend=biber,
  style=verbose-note,    % 脚注内にフル書誌→短縮形（人文系標準）
  autocite=footnote,
  defernumbers=true,     % 分割出力時の番号管理
  sorting=nyt,
  abbreviate=false,
  language=auto,         % langid による言語別自動切替
]{biblatex}

% 参考文献一覧の見出し階層定義
\defbibheading{subbibliography}[\bibname]{%
  \section*{#1}}
\defbibheading{subsubbibliography}[\bibname]{%
  \subsection*{#1}}

%%% 書誌フォーマット（修論準拠）%%%
%
% langid による言語別自動切替:
%   french  → « guillemets », "dans", "et", "éd."
%   english → "quotation marks", "in", "and", "ed."
%   japanese → キーワードで別処理
% .bibエントリに langid = {french/english/japanese} を指定すること。

% 脚注・参考文献一覧ともに姓-名の順に統一
\DeclareNameAlias{sortname}{family-given}
\DeclareNameAlias{default}{family-given}

% 欧語: 姓をスモールキャピタルにする
% \ifkeyword は短縮引用で効かないため、LuaLaTeX で文字コード判定
\directlua{
  function is_cjk(s)
    local first = utf8.codepoint(s, 1)
    if first >= 0x3000 and first <= 0x9FFF then return true end
    if first >= 0xF900 and first <= 0xFAFF then return true end
    return false
  end
}
\renewcommand*{\mkbibnamefamily}[1]{%
  \directlua{
    if is_cjk("\luaescapestring{#1}") then
      tex.sprint("\luaescapestring{#1}")
    else
      tex.sprint("\\textsc{\luaescapestring{#1}}")
    end
  }%
}

% 日本語文献: 姓名スペース除去 + 句読点を日本語に
% 日本語用の句読点設定マクロ
\newcommand*{\setjapanesepunct}{%
  \renewcommand*{\revsdnamepunct}{}%
  \renewcommand*{\bibnamedelimd}{}%
  \renewcommand*{\newunitpunct}{\addspace}%
  \renewcommand*{\labelnamepunct}{\addspace}%
  \renewcommand*{\subtitlepunct}{：}%
  \renewcommand*{\bibpagespunct}{、}%
  \renewcommand*{\multicitedelim}{、}%
  \renewcommand*{\finentrypunct}{。}%
  \DeclareDelimFormat{nameyeardelim}{、}%
  \DeclareDelimFormat{nonameyeardelim}{、}%
  \DeclareDelimFormat[bib,biblist]{multinamedelim}{・}%
  \DeclareDelimFormat[bib,biblist]{finalnamedelim}{・}%
  \DeclareDelimFormat{titleyeardelim}{、}%
  \DeclareDelimFormat{publocpunct}{、}%
  \DeclareDelimFormat{publocdelim}{、}%
  \DeclareDelimFormat{publisherdelim}{、}%
  \renewcommand*{\newblockpunct}{\addspace}%
  \DeclareDelimFormat{postnotedelim}{、}%
}
% 欧語用の句読点設定マクロ
\newcommand*{\setwesternpunct}{%
  \renewcommand*{\revsdnamepunct}{\addcomma\space}%
  \renewcommand*{\bibnamedelimd}{\addspace}%
  \renewcommand*{\newunitpunct}{\adddot\space}%
  \renewcommand*{\subtitlepunct}{\addcolon\space}%
  \renewcommand*{\finentrypunct}{\adddot}%
}
% 日本語: 出版年の後に「年」を付ける
\DeclareFieldFormat{datetext}{%
  \ifkeyword{japanese}{#1年}{#1}}
\renewbibmacro*{date}{%
  \printtext[datetext]{\printdate}%
}

\AtEveryBibitem{\ifkeyword{japanese}{\setjapanesepunct}{\setwesternpunct}}
\AtEveryCitekey{\ifkeyword{japanese}{\setjapanesepunct}{\setwesternpunct}}

% 脚注末尾: 日本語「。」/ 欧語「.」
\renewcommand{\bibfootnotewrapper}[1]{\bibsentence #1}

% 短縮引用: 日本語では「著者、前掲書」形式にする
\renewbibmacro*{cite:short}{%
  \ifkeyword{japanese}%
    {\printnames{labelname}、前掲書}%
    {\printnames{labelname}%
     \setunit*{\printdelim{nametitledelim}}%
     \printtext[bibhyperref]{%
       \printfield[citetitle]{labeltitle}}}%
}

% autocite: 末尾句点を postnote の後に出す
\DeclareCiteCommand{\autocite}[\mkbibfootnote]
  {\usebibmacro{prenote}}
  {\usebibmacro{citeindex}%
   \usebibmacro{cite}}
  {\multicitedelim}
  {\usebibmacro{cite:postnote}\finentrypunct}

% 出版社と年の間の区切りを日本語では「、」に
\renewbibmacro*{publisher+location+date}{%
  \printlist{location}%
  \iflistundef{publisher}%
    {\setunit*{\addcomma\space}}%
    {\setunit*{\ifkeyword{japanese}{、}{\addcomma\space}}}%
  \printlist{publisher}%
  \setunit*{\ifkeyword{japanese}{、}{\addcomma\space}}%
  \usebibmacro{date}%
  \newunit%
}

% 論文・章題の引用符
% 日本語「」、フランス語 « »、英語 " "
\newcommand*{\langquote}[1]{%
  \iffieldequalstr{langid}{french}%
    {\guillemotleft\addnbspace #1\addnbspace\guillemotright}%
    {\mkbibquote{#1}}%
}
\DeclareFieldFormat[article]{title}{%
  \ifkeyword{japanese}{「#1」}{\langquote{#1}}}
\DeclareFieldFormat[incollection]{title}{%
  \ifkeyword{japanese}{「#1」}{\langquote{#1}}}
\DeclareFieldFormat[inproceedings]{title}{%
  \ifkeyword{japanese}{「#1」}{\langquote{#1}}}
\DeclareFieldFormat[thesis]{title}{%
  \ifkeyword{japanese}{「#1」}{\langquote{#1}}}

% 日本語の書名を『』で囲む（イタリック→ゴシック化を防止）
\DeclareFieldFormat[book]{title}{%
  \ifkeyword{japanese}{『#1』}{\mkbibemph{#1}}}
\DeclareFieldFormat[book]{subtitle}{%
  \ifkeyword{japanese}{#1}{\mkbibemph{#1}}}
\DeclareFieldFormat[incollection]{booktitle}{%
  \ifkeyword{japanese}{『#1』}{\mkbibemph{#1}}}
\DeclareFieldFormat[book]{citetitle}{%
  \ifkeyword{japanese}{『#1』}{\mkbibemph{#1}}}
\DeclareFieldFormat[incollection]{citetitle}{%
  \ifkeyword{japanese}{「#1」}{\langquote{#1}}}
\DeclareFieldFormat[article]{citetitle}{%
  \ifkeyword{japanese}{「#1」}{\langquote{#1}}}

% 日本語: "In:" を抑制、欧語: 大文字 "In:" を維持
% 日本語: 抑制、フランス語: dans、英語: in
\renewbibmacro*{in:}{%
  \ifkeyword{japanese}%
    {}%
    {\iffieldequalstr{langid}{french}%
      {\printtext{dans\intitlepunct}}%
      {\printtext{in\intitlepunct}}}%
}

% 日本語: 短縮引用「前掲注X」+ pages 表記
\DefineBibliographyStrings{english}{%
  seenote = {前掲注},
}
% 日本語文献の pages を日本語化
\DeclareFieldFormat[incollection]{pages}{%
  \ifkeyword{japanese}{#1頁}{pp.\addnbspace #1}}
\DeclareFieldFormat[article]{pages}{%
  \ifkeyword{japanese}{#1頁}{pp.\addnbspace #1}}

% 日本語: "Edited by ..." を抑制、欧語: 名前 (eds.) 形式
\renewbibmacro*{byeditor+others}{%
  \ifkeyword{japanese}%
    {}%
    {\ifnameundef{editor}%
      {}%
      {\printnames[byeditor]{editor}%
       \addspace\printeditorlabel
       \clearname{editor}%
       \newunit}%
    }%
}

% 編者の括弧付き表記ヘルパー
\newcommand*{\printeditorlabel}{%
  \iffieldequalstr{langid}{french}%
    {\printtext[parens]{\ifnumgreater{\value{editor}}{1}{éds\adddot}{éd\adddot}}}%
    {\printtext[parens]{\ifnumgreater{\value{editor}}{1}{eds\adddot}{ed\adddot}}}%
}

% editor+othersstrg / editor+others: 括弧付き (eds.)/(éds.) を出力
\renewbibmacro*{editor+othersstrg}{\printeditorlabel}
\renewbibmacro*{editor+others}{%
  \ifboolexpr{
    test \ifuseeditor
    and
    not test {\ifnameundef{editor}}
  }
    {\printnames{editor}%
     \addspace\printeditorlabel
     \clearname{editor}}
    {}%
}

% book型: 編者名 (éds.)/(eds.) の形式
\renewbibmacro*{bbx:editor}{%
  \ifnameundef{editor}{}%
    {\printnames{editor}%
     \usebibmacro{editor+othersstrg}%
     \clearname{editor}%
     \setunit{\labelnamepunct}}%
}

% フランス語文献: "and" → "et"
\AtEveryBibitem{%
  \iffieldequalstr{langid}{french}%
    {\DeclareDelimFormat{finalnamedelim}{\addspace et\addspace}}%
    {}%
}
\AtEveryCitekey{%
  \iffieldequalstr{langid}{french}%
    {\DeclareDelimFormat{finalnamedelim}{\addspace et\addspace}}%
    {}%
}

%%% 組版 %%%
\usepackage{setspace}
\onehalfspacing
\usepackage[
  top=30mm,bottom=30mm,
  left=30mm,right=25mm,
]{geometry}


%%% 引用符 %%%
\usepackage[autostyle=true, french=guillemets, english=american]{csquotes}

%%% ハイパーリンク %%%
\usepackage[dvipsnames,svgnames,x11names]{xcolor}
\usepackage{hyperref}
\hypersetup{
  colorlinks=true,
  linkcolor=black,
  citecolor=black,
  urlcolor=blue!60!black,
  bookmarksnumbered=true,
}

%%% その他 %%%
\usepackage{graphicx}
\usepackage{booktabs}
\usepackage{wrapfig}

%%% 見出し %%%
\setcounter{secnumdepth}{1}
\setcounter{tocdepth}{2}
  ← 各プロジェクトの main.tex から

%%% フォント %%%
\usepackage{luatexja-fontspec}
\usepackage{iftex}
% 和文明朝: 原ノ味明朝（TeX Live同梱、全OS対応）
\setmainjfont{HaranoAjiMincho}
% 和文ゴシック: OS別にフォールバック
%   macOS: ヒラギノ角ゴシック ProN
%   Linux/Windows: Noto Sans CJK JP → 原ノ味ゴシック（TeX Live同梱）
\IfFontExistsTF{Hiragino Kaku Gothic ProN}{%
  \setsansjfont{Hiragino Kaku Gothic ProN}[
    UprightFont = {Hiragino Kaku Gothic ProN W3},
    BoldFont    = {Hiragino Kaku Gothic ProN W6},
  ]%
}{%
  \IfFontExistsTF{Noto Sans CJK JP}{%
    \setsansjfont{Noto Sans CJK JP}%
  }{%
    \setsansjfont{HaranoAjiGothic}%
  }%
}
% 見出しをゴシック体にする
\AddToHook{begindocument/before}{%
  \ModifyHeading{chapter}{font={\sffamily\bfseries}}%
  \ModifyHeading{section}{font={\sffamily\bfseries}}%
  \ModifyHeading{subsection}{font={\sffamily\bfseries}}%
}

% 部の見出し: 番号と題名を別行に、行間を広げる
\makeatletter
\renewcommand{\part}{%
  \clearpage
  \thispagestyle{empty}%
  \null\vfil
  \secdef\@part\@spart}
\def\@part[#1]#2{%
  \refstepcounter{part}%
  \addcontentsline{toc}{part}{\thepart\hspace{1em}#1}%
  \begin{center}
    {\huge\sffamily\bfseries 第\thepart 部}\\[1.5em]
    {\LARGE\sffamily\bfseries #2}
  \end{center}
  \vfil}
\def\@spart#1{%
  \begin{center}
    {\LARGE\sffamily\bfseries #1}
  \end{center}
  \vfil}
\makeatother
% 欧文: TeX Gyre Termes（Times系、TeX Live同梱）
\setmainfont{TeX Gyre Termes}

%%% 脚注番号の括弧を外す %%%
\renewcommand{\thefootnote}{\arabic{footnote}}

%%% 節見出しの番号を「第n節」にする %%%
\renewcommand{\thesection}{第\arabic{section}節}
\renewcommand{\thesubsection}{\arabic{section}.\arabic{subsection}}

%%% 目次の体裁 %%%
% 部: 章と同じインデント体裁で、太字ゴシック
\makeatletter
\renewcommand*{\l@part}[2]{%
  \ifnum \c@tocdepth >-2\relax
    \addpenalty{-\@highpenalty}%
    \addvspace{1em \@plus\p@}%
    \@dottedtocline{-1}{0em}{5em}{\bfseries\sffamily #1}{\bfseries #2}%
    \nobreak
  \fi}
\makeatother

% oneside: 空白ページの挿入を防止
\let\cleardoublepage\clearpage

%%% 参考文献 %%%
\usepackage[
  backend=biber,
  style=verbose-note,    % 脚注内にフル書誌→短縮形（人文系標準）
  autocite=footnote,
  defernumbers=true,     % 分割出力時の番号管理
  sorting=nyt,
  abbreviate=false,
  language=auto,         % langid による言語別自動切替
]{biblatex}

% 参考文献一覧の見出し階層定義
\defbibheading{subbibliography}[\bibname]{%
  \section*{#1}}
\defbibheading{subsubbibliography}[\bibname]{%
  \subsection*{#1}}

%%% 書誌フォーマット（修論準拠）%%%
%
% langid による言語別自動切替:
%   french  → « guillemets », "dans", "et", "éd."
%   english → "quotation marks", "in", "and", "ed."
%   japanese → キーワードで別処理
% .bibエントリに langid = {french/english/japanese} を指定すること。

% 脚注・参考文献一覧ともに姓-名の順に統一
\DeclareNameAlias{sortname}{family-given}
\DeclareNameAlias{default}{family-given}

% 欧語: 姓をスモールキャピタルにする
% \ifkeyword は短縮引用で効かないため、LuaLaTeX で文字コード判定
\directlua{
  function is_cjk(s)
    local first = utf8.codepoint(s, 1)
    if first >= 0x3000 and first <= 0x9FFF then return true end
    if first >= 0xF900 and first <= 0xFAFF then return true end
    return false
  end
}
\renewcommand*{\mkbibnamefamily}[1]{%
  \directlua{
    if is_cjk("\luaescapestring{#1}") then
      tex.sprint("\luaescapestring{#1}")
    else
      tex.sprint("\\textsc{\luaescapestring{#1}}")
    end
  }%
}

% 日本語文献: 姓名スペース除去 + 句読点を日本語に
% 日本語用の句読点設定マクロ
\newcommand*{\setjapanesepunct}{%
  \renewcommand*{\revsdnamepunct}{}%
  \renewcommand*{\bibnamedelimd}{}%
  \renewcommand*{\newunitpunct}{\addspace}%
  \renewcommand*{\labelnamepunct}{\addspace}%
  \renewcommand*{\subtitlepunct}{：}%
  \renewcommand*{\bibpagespunct}{、}%
  \renewcommand*{\multicitedelim}{、}%
  \renewcommand*{\finentrypunct}{。}%
  \DeclareDelimFormat{nameyeardelim}{、}%
  \DeclareDelimFormat{nonameyeardelim}{、}%
  \DeclareDelimFormat[bib,biblist]{multinamedelim}{・}%
  \DeclareDelimFormat[bib,biblist]{finalnamedelim}{・}%
  \DeclareDelimFormat{titleyeardelim}{、}%
  \DeclareDelimFormat{publocpunct}{、}%
  \DeclareDelimFormat{publocdelim}{、}%
  \DeclareDelimFormat{publisherdelim}{、}%
  \renewcommand*{\newblockpunct}{\addspace}%
  \DeclareDelimFormat{postnotedelim}{、}%
}
% 欧語用の句読点設定マクロ
\newcommand*{\setwesternpunct}{%
  \renewcommand*{\revsdnamepunct}{\addcomma\space}%
  \renewcommand*{\bibnamedelimd}{\addspace}%
  \renewcommand*{\newunitpunct}{\adddot\space}%
  \renewcommand*{\subtitlepunct}{\addcolon\space}%
  \renewcommand*{\finentrypunct}{\adddot}%
}
% 日本語: 出版年の後に「年」を付ける
\DeclareFieldFormat{datetext}{%
  \ifkeyword{japanese}{#1年}{#1}}
\renewbibmacro*{date}{%
  \printtext[datetext]{\printdate}%
}

\AtEveryBibitem{\ifkeyword{japanese}{\setjapanesepunct}{\setwesternpunct}}
\AtEveryCitekey{\ifkeyword{japanese}{\setjapanesepunct}{\setwesternpunct}}

% 脚注末尾: 日本語「。」/ 欧語「.」
\renewcommand{\bibfootnotewrapper}[1]{\bibsentence #1}

% 短縮引用: 日本語では「著者、前掲書」形式にする
\renewbibmacro*{cite:short}{%
  \ifkeyword{japanese}%
    {\printnames{labelname}、前掲書}%
    {\printnames{labelname}%
     \setunit*{\printdelim{nametitledelim}}%
     \printtext[bibhyperref]{%
       \printfield[citetitle]{labeltitle}}}%
}

% autocite: 末尾句点を postnote の後に出す
\DeclareCiteCommand{\autocite}[\mkbibfootnote]
  {\usebibmacro{prenote}}
  {\usebibmacro{citeindex}%
   \usebibmacro{cite}}
  {\multicitedelim}
  {\usebibmacro{cite:postnote}\finentrypunct}

% 出版社と年の間の区切りを日本語では「、」に
\renewbibmacro*{publisher+location+date}{%
  \printlist{location}%
  \iflistundef{publisher}%
    {\setunit*{\addcomma\space}}%
    {\setunit*{\ifkeyword{japanese}{、}{\addcomma\space}}}%
  \printlist{publisher}%
  \setunit*{\ifkeyword{japanese}{、}{\addcomma\space}}%
  \usebibmacro{date}%
  \newunit%
}

% 論文・章題の引用符
% 日本語「」、フランス語 « »、英語 " "
\newcommand*{\langquote}[1]{%
  \iffieldequalstr{langid}{french}%
    {\guillemotleft\addnbspace #1\addnbspace\guillemotright}%
    {\mkbibquote{#1}}%
}
\DeclareFieldFormat[article]{title}{%
  \ifkeyword{japanese}{「#1」}{\langquote{#1}}}
\DeclareFieldFormat[incollection]{title}{%
  \ifkeyword{japanese}{「#1」}{\langquote{#1}}}
\DeclareFieldFormat[inproceedings]{title}{%
  \ifkeyword{japanese}{「#1」}{\langquote{#1}}}
\DeclareFieldFormat[thesis]{title}{%
  \ifkeyword{japanese}{「#1」}{\langquote{#1}}}

% 日本語の書名を『』で囲む（イタリック→ゴシック化を防止）
\DeclareFieldFormat[book]{title}{%
  \ifkeyword{japanese}{『#1』}{\mkbibemph{#1}}}
\DeclareFieldFormat[book]{subtitle}{%
  \ifkeyword{japanese}{#1}{\mkbibemph{#1}}}
\DeclareFieldFormat[incollection]{booktitle}{%
  \ifkeyword{japanese}{『#1』}{\mkbibemph{#1}}}
\DeclareFieldFormat[book]{citetitle}{%
  \ifkeyword{japanese}{『#1』}{\mkbibemph{#1}}}
\DeclareFieldFormat[incollection]{citetitle}{%
  \ifkeyword{japanese}{「#1」}{\langquote{#1}}}
\DeclareFieldFormat[article]{citetitle}{%
  \ifkeyword{japanese}{「#1」}{\langquote{#1}}}

% 日本語: "In:" を抑制、欧語: 大文字 "In:" を維持
% 日本語: 抑制、フランス語: dans、英語: in
\renewbibmacro*{in:}{%
  \ifkeyword{japanese}%
    {}%
    {\iffieldequalstr{langid}{french}%
      {\printtext{dans\intitlepunct}}%
      {\printtext{in\intitlepunct}}}%
}

% 日本語: 短縮引用「前掲注X」+ pages 表記
\DefineBibliographyStrings{english}{%
  seenote = {前掲注},
}
% 日本語文献の pages を日本語化
\DeclareFieldFormat[incollection]{pages}{%
  \ifkeyword{japanese}{#1頁}{pp.\addnbspace #1}}
\DeclareFieldFormat[article]{pages}{%
  \ifkeyword{japanese}{#1頁}{pp.\addnbspace #1}}

% 日本語: "Edited by ..." を抑制、欧語: 名前 (eds.) 形式
\renewbibmacro*{byeditor+others}{%
  \ifkeyword{japanese}%
    {}%
    {\ifnameundef{editor}%
      {}%
      {\printnames[byeditor]{editor}%
       \addspace\printeditorlabel
       \clearname{editor}%
       \newunit}%
    }%
}

% 編者の括弧付き表記ヘルパー
\newcommand*{\printeditorlabel}{%
  \iffieldequalstr{langid}{french}%
    {\printtext[parens]{\ifnumgreater{\value{editor}}{1}{éds\adddot}{éd\adddot}}}%
    {\printtext[parens]{\ifnumgreater{\value{editor}}{1}{eds\adddot}{ed\adddot}}}%
}

% editor+othersstrg / editor+others: 括弧付き (eds.)/(éds.) を出力
\renewbibmacro*{editor+othersstrg}{\printeditorlabel}
\renewbibmacro*{editor+others}{%
  \ifboolexpr{
    test \ifuseeditor
    and
    not test {\ifnameundef{editor}}
  }
    {\printnames{editor}%
     \addspace\printeditorlabel
     \clearname{editor}}
    {}%
}

% book型: 編者名 (éds.)/(eds.) の形式
\renewbibmacro*{bbx:editor}{%
  \ifnameundef{editor}{}%
    {\printnames{editor}%
     \usebibmacro{editor+othersstrg}%
     \clearname{editor}%
     \setunit{\labelnamepunct}}%
}

% フランス語文献: "and" → "et"
\AtEveryBibitem{%
  \iffieldequalstr{langid}{french}%
    {\DeclareDelimFormat{finalnamedelim}{\addspace et\addspace}}%
    {}%
}
\AtEveryCitekey{%
  \iffieldequalstr{langid}{french}%
    {\DeclareDelimFormat{finalnamedelim}{\addspace et\addspace}}%
    {}%
}

%%% 組版 %%%
\usepackage{setspace}
\onehalfspacing
\usepackage[
  top=30mm,bottom=30mm,
  left=30mm,right=25mm,
]{geometry}


%%% 引用符 %%%
\usepackage[autostyle=true, french=guillemets, english=american]{csquotes}

%%% ハイパーリンク %%%
\usepackage[dvipsnames,svgnames,x11names]{xcolor}
\usepackage{hyperref}
\hypersetup{
  colorlinks=true,
  linkcolor=black,
  citecolor=black,
  urlcolor=blue!60!black,
  bookmarksnumbered=true,
}

%%% その他 %%%
\usepackage{graphicx}
\usepackage{booktabs}
\usepackage{wrapfig}

%%% 見出し %%%
\setcounter{secnumdepth}{1}
\setcounter{tocdepth}{2}
